\section{Draft: Creating a Program to read MNIST numbers}

First, we initialize an MLP with an input size and layer params. The layer params are 
simply the size of each layer. 

After that, we run a training loop. A training loop is simply composed of these fundamental steps:

sequentual bullet points.

Zero out all of the local gradients. 
Feedforward the input through the MLP an obtain an output.
Calculate the loss with the output.
Potentially reduce the loss with gradient descent.

Repeat with a new batch of data.

Time to break down each indidivual step.
We must zero out all of the local gradients because we don't want any past calculations to 
interfere. Basically just a clean slate.

Feed forward is just running the input through the MLP. This is usually done in matrices or tensors.

With the output, we can calculate how far off we are and grade the model. Our goal for the MLP is to 
reduce the loss calculated. We can punish the model by having a larger loss. This brings us to loss functions
that use this relationship of the magnitude of wrongness to the magnitude of loss. There is simple error loss,
mean squared loss, and cross entropy loss. 

Now we must reduce the loss as much as possible which is in the direction of the gradients. We run back 
propogration through each operaiton to obtain all the local derivatives. Think of the local derviatives as 
a ratio of how much each individual component helps toward the output. Almost like ingredients toward a dish.
After, simply adjust the weights and biases in the direction of the gradient. If we want to increase the value,
positive gradient and decrease the value, negative gradient. Comes in theory with the directional derivative and 
the dot product that has a cosine function. 

\subsection{MNIST training loop}

In the case with the NNIST loop, we go through each component of the training loop one by one and build it.
When buildinging, it's important to test each component instead of writing all at once. First we obtain the data
as a 784 (28 x 28 image) vector and additionally normalize every value which is to divide by 255. We get this number 
because the data is in the form of brightness values from 0-255 hence the normalization factor. We additionally obtain 
the proper correct label to check in the loss function. 

Now feed forward and obtain an output vector. I decided for an 10 size output vector for each number from 0-9. We run this
softmax to gain a "percentage" distrubution of numbers and to clamp them from 0-1. This is import due to the one-hot labelling
used for calculating the loss.

I implemented both mean squared and cross entropy loss. In comparison with the one-hot label, we can compare that with our 10 size
output vector. Mean squared simply sums all of the the squared difference between each index and normalizing it with the size. 
(Hence the name mean square.) Cross entropy takes the sum of the expected value multiplied by the log of the predicted output value.

Given here are the algorithms.

\begin{definition}[Mean Square Loss]
    \begin{equation}
        \mathcal{L}(x) = \frac{\sum_{i}^{n} (\text{expected}_i - \text{predicted}_i)^{2}}{n} 
    \end{equation}
\end{definition}

\begin{definition}[Cross Entropy Loss]
    \begin{equation}
        \mathcal{L}(x) = \frac{\sum_{i}^{n} (\text{expected}_i \cdot \log\left(\text{softmax}(\text{predicted}_i) \right))^{2}}{n} 
    \end{equation}
\end{definition}

Run back propogration and now gradient descent.
\begin{align*}
    D_{\mathbf{\hat{u}}} \, f  =& \nabla f \cdot \hat{u} \\
                                  =& |\nabla f| \, \cancelto{1}{|\hat{u}|} \cos\left(\theta\right), \quad -1 \leq \cos\left(\theta\right) f \leq 1
\end{align*}

\begin{equation}
    \Longrightarrow - \nabla f \leq D_{\mathbf{\hat{u}}} \, f \leq \nabla f
\end{equation}

